% GENERATED FILE. Edit canonical lessons, then run npm run generate:books.
% canonical-source-hash: fnv1a64:ecf6a733fed44d83
% canonical-lessons: PT-C107-sandes, PT-C107-molho, PT-C107-vinagre, PT-C107-sobremesa, PT-C107-refeicao

\chapter{Sandwich, Sauce, Vinegar, Dessert, Meal}
\label{ch:pt-sandwich-sauce-vinegar-dessert-meal}
\hlchaptermodality{107}

\begin{chapteropening}
\textbf{By the end of this chapter:} \emph{I can say a sandwich, a sauce, vinegar, a dessert and a meal.}

Complete the last lesson of chapter 107: I can say a sandwich, a sauce, vinegar, a dessert and a meal.
\end{chapteropening}

\section[a sandes]{a sandes — a sandwich}

\label{lesson:PT-C107-sandes}



\begin{quote}
Before the new one: say the Portuguese for chocolate, then the Portuguese for a biscuit.
\end{quote}

\subsection*{You'll want to know: a sandes}
\textbf{a sandes} — \textquotedblleft{}a sandwich\textquotedblright{}. In Brazil, \textbf{o sanduíche}.

\begin{cousinweb}
\textbf{sandes} is a Portuguese reshaping of English \emph{sandwich}, named after the Earl of Sandwich — borrowed.
\end{cousinweb}

\begin{grammarlens}[title={What you've built}]
One more word for this chapter.
\end{grammarlens}

\subsection*{Your turn}
\begin{itemize}
\raggedright
  \item \emph{Say it:} \emph{a sandes}
  \item \emph{Say it:} \emph{a sandes}, once more
  \item \emph{Say it:} say \emph{a sandes}
  \item \emph{Recall:} say the Portuguese for chocolate, then the Portuguese for a biscuit, then say \emph{a sandes} again
\end{itemize}

\begin{tcolorbox}[breakable,colback=teal!4,colframe=teal!35!black,title={Before you move on}]
What does a sandes mean? (\textquotedblleft{}A sandwich\textquotedblright{}.) Say it once more.
\end{tcolorbox}
\section[o molho]{o molho — a sauce}

\label{lesson:PT-C107-molho}



\begin{quote}
Before the new one: say the Portuguese for a biscuit, then the Portuguese for a sandwich.
\end{quote}

\subsection*{You'll want to know: o molho}
\textbf{o molho} — \textquotedblleft{}a sauce\textquotedblright{}.

\begin{cousinweb}
\textbf{molho} is built from the verb \textquotedblleft{}to wet, to soak\textquotedblright{} (EtymDB 2.1), which $\leftarrow$ Vulgar Latin \emph{*molliare}, \textquotedblleft{}to soften by soaking\textquotedblright{}, from Latin \emph{mollis}, \textquotedblleft{}soft\textquotedblright{} (EtymDB 2.1).
\end{cousinweb}

\begin{grammarlens}[title={What you've built}]
One more word for this chapter.
\end{grammarlens}

\subsection*{Your turn}
\begin{itemize}
\raggedright
  \item \emph{Say it:} \emph{o molho}
  \item \emph{Say it:} \emph{o molho}, once more
  \item \emph{Say it:} say \emph{o molho}
  \item \emph{Recall:} say the Portuguese for a biscuit, then the Portuguese for a sandwich, then say \emph{o molho} again
\end{itemize}

\begin{tcolorbox}[breakable,colback=teal!4,colframe=teal!35!black,title={Before you move on}]
What does o molho mean? (\textquotedblleft{}A sauce\textquotedblright{}.) Say it once more.
\end{tcolorbox}
\section[o vinagre]{o vinagre — vinegar}

\label{lesson:PT-C107-vinagre}



\begin{quote}
Before the new one: say the Portuguese for a sandwich, then the Portuguese for a sauce.
\end{quote}

\subsection*{You'll want to know: o vinagre}
\textbf{o vinagre} — \textquotedblleft{}vinegar\textquotedblright{}.

\begin{cousinweb}
\textbf{vinagre} $\leftarrow$ Catalan \emph{vinagre}, \textquotedblleft{}sour wine\textquotedblright{} — borrowed, from Latin \emph{vinum}, \textquotedblleft{}wine\textquotedblright{}, and \emph{acer}, \textquotedblleft{}sour\textquotedblright{}. EtymDB 2.1 traces it only to Old Portuguese \emph{vinagre}.
\end{cousinweb}

\begin{grammarlens}[title={What you've built}]
One more word for this chapter.
\end{grammarlens}

\subsection*{Your turn}
\begin{itemize}
\raggedright
  \item \emph{Say it:} \emph{o vinagre}
  \item \emph{Say it:} \emph{o vinagre}, once more
  \item \emph{Say it:} say \emph{o vinagre}
  \item \emph{Recall:} say the Portuguese for a sandwich, then the Portuguese for a sauce, then say \emph{o vinagre} again
\end{itemize}

\begin{tcolorbox}[breakable,colback=teal!4,colframe=teal!35!black,title={Before you move on}]
What does o vinagre mean? (\textquotedblleft{}Vinegar\textquotedblright{}.) Say it once more.
\end{tcolorbox}
\section[a sobremesa]{a sobremesa — a dessert}

\label{lesson:PT-C107-sobremesa}



\begin{quote}
Before the new one: say the Portuguese for a sauce, then the Portuguese for vinegar.
\end{quote}

\subsection*{You'll want to know: a sobremesa}
\textbf{a sobremesa} — \textquotedblleft{}a dessert\textquotedblright{}.

\begin{cousinweb}
\textbf{sobremesa} is \textquotedblleft{}over the table\textquotedblright{}: built from \emph{sobre}, \textquotedblleft{}over\textquotedblright{} ($\leftarrow$ Latin \emph{super}), and \emph{mesa}, \textquotedblleft{}table\textquotedblright{} — what comes when the meal is done.
\end{cousinweb}

\begin{grammarlens}[title={What you've built}]
One more word for this chapter.
\end{grammarlens}

\subsection*{Your turn}
\begin{itemize}
\raggedright
  \item \emph{Say it:} \emph{a sobremesa}
  \item \emph{Say it:} \emph{a sobremesa}, once more
  \item \emph{Say it:} say \emph{a sobremesa}
  \item \emph{Recall:} say the Portuguese for a sauce, then the Portuguese for vinegar, then say \emph{a sobremesa} again
\end{itemize}

\begin{tcolorbox}[breakable,colback=teal!4,colframe=teal!35!black,title={Before you move on}]
What does a sobremesa mean? (\textquotedblleft{}A dessert\textquotedblright{}.) Say it once more.
\end{tcolorbox}
\section[a refeição]{a refeição — a meal}

\label{lesson:PT-C107-refeicao}



\begin{quote}
Before the new one: say the Portuguese for vinegar, then the Portuguese for a dessert.
\end{quote}

\subsection*{You'll want to know: a refeição}
\textbf{a refeição} — \textquotedblleft{}a meal\textquotedblright{}.

\begin{cousinweb}
\textbf{refeição} $\leftarrow$ Latin \emph{refectio, refectionis}, \textquotedblleft{}restoring\textquotedblright{} (EtymDB 2.1), from \emph{reficere}, \textquotedblleft{}to remake\textquotedblright{} — a learned borrowing from Latin.
\end{cousinweb}

\begin{grammarlens}[title={What you've built}]
One more word for this chapter.
\end{grammarlens}

\subsection*{Your turn}
\begin{itemize}
\raggedright
  \item \emph{Say it:} \emph{a refeição}
  \item \emph{Say it:} \emph{a refeição}, once more
  \item \emph{Say it:} say \emph{a refeição}
  \item \emph{Recall:} say the Portuguese for vinegar, then the Portuguese for a dessert, then say \emph{a refeição} again
\end{itemize}

\begin{tcolorbox}[breakable,colback=teal!4,colframe=teal!35!black,title={Before you move on}]
What does a refeição mean? (\textquotedblleft{}A meal\textquotedblright{}.) Say it once more.
\end{tcolorbox}
